\textbf{Phần 1. Trắc nghiệm lựa chọn}
\setcounterpageref{ex}{0}

\begin{ex} %[3D2V36]
Tính giá trị của biểu thức: $A = 125 \times 4 - 180$.
\choice
{$300$}
{\True $320$}
{$350$}
{$280$}
\loigiai{
$125 \times 4 = 500$; $500 - 180 = 320$.
\textbf{Đáp án của câu hỏi là B}
}
\end{ex}

\begin{ex} %[3D2V36]
Một chuyến xe buýt chở được nhiều nhất 45 hành khách. Hỏi 6 chuyến xe buýt như thế chở được nhiều nhất bao nhiêu hành khách?
\choice
{$240$ hành khách}
{$250$ hành khách}
{$280$ hành khách}
{\True $270$ hành khách}
\loigiai{
Số hành khách nhiều nhất chở được là: $45 \times 6 = 270$ (hành khách).
\textbf{Đáp án của câu hỏi là D}
}
\end{ex}

\begin{ex} %[3D2V36]
Tìm $x$, biết: $x : 3 = 142 + 28$.
\choice
{\True $510$}
{$480$}
{$520$}
{$500$}
\loigiai{
$x : 3 = 170 \Rightarrow x = 170 \times 3 = 510$.
\textbf{Đáp án của câu hỏi là A}
}
\end{ex}

\begin{ex} %[3D2V36]
Một cuộn dây cáp dài $135\text{ m}$. Người ta dùng 4 cuộn cáp như thế để lắp đường dây điện. Chiều dài đường dây điện đã lắp là
\choice
{$500\text{ m}$}
{$520\text{ m}$}
{\True $540\text{ m}$}
{$560\text{ m}$}
\loigiai{
Tổng chiều dài cáp: $135 \times 4 = 540\text{ (m)}$.
\textbf{Đáp án của câu hỏi là C}
}
\end{ex}

\begin{ex} %[3D2V36]
Một hình vuông có cạnh dài $145\text{ cm}$. Chu vi của hình vuông đó là
\choice
{$560\text{ cm}$}
{\True $580\text{ cm}$}
{$600\text{ cm}$}
{$540\text{ cm}$}
\loigiai{
Chu vi hình vuông: $145 \times 4 = 580\text{ (cm)}$.
\textbf{Đáp án của câu hỏi là B}
}
\end{ex}

\begin{ex} %[3D2V36]
Từ bốn chữ số $1, 2, 3, 4$, người ta lập một số có ba chữ số và một số có một chữ số để được tích lớn nhất có thể. Tích lớn nhất đó là
\choice
{$964$}
{$1248$}
{$1324$}
{\True $1284$}
\loigiai{
Chọn thừa số một chữ số lớn nhất là $4$, số có 3 chữ số lớn nhất tạo từ $1, 2, 3$ là $321$. Tích lớn nhất: $321 \times 4 = 1284$.
\textbf{Đáp án của câu hỏi là D}
}
\end{ex}

\begin{ex} %[3D2V36]
Một cửa hàng nhập về 3 kiện hàng, mỗi kiện có 150 hộp sữa. Cửa hàng đã bán được 280 hộp sữa. Cửa hàng còn lại bao nhiêu hộp sữa?
\choice
{\True $170$ hộp}
{$150$ hộp}
{$180$ hộp}
{$190$ hộp}
\loigiai{
Tổng số hộp sữa: $150 \times 3 = 450$ hộp. Còn lại: $450 - 280 = 170$ hộp sữa.
\textbf{Đáp án của câu hỏi là A}
}
\end{ex}

\begin{ex} %[3D2V36]
Tính nhanh giá trị của biểu thức: $115 \times 2 + 115 \times 3 + 115$.
\choice
{$575$}
{$650$}
{\True $690$}
{$720$}
\loigiai{
$115 \times (2 + 3 + 1) = 115 \times 6 = 690$.
\textbf{Đáp án của câu hỏi là C}
}
\end{ex}

\begin{ex} %[3D2V36]
Mỗi bao gạo cân nặng $50\text{ kg}$. Một xe tải chở 8 bao gạo và 4 bao ngô, mỗi bao ngô nặng $40\text{ kg}$. Tổng khối lượng hàng trên xe tải là
\choice
{$480\text{ kg}$}
{\True $560\text{ kg}$}
{$520\text{ kg}$}
{$600\text{ kg}$}
\loigiai{
Khối lượng gạo: $50 \times 8 = 400\text{ kg}$. Khối lượng ngô: $40 \times 4 = 160\text{ kg}$. Tổng khối lượng: $400 + 160 = 560\text{ (kg)}$.
\textbf{Đáp án của câu hỏi là B}
}
\end{ex}

\begin{ex} %[3D2V36]
Nếu một thừa số tăng thêm 3 đơn vị và thừa số kia giữ nguyên bằng 124 thì tích tăng thêm bao nhiêu đơn vị?
\choice
{$124$}
{$248$}
{$496$}
{\True $372$}
\loigiai{
Tích tăng thêm bằng: $124 \times 3 = 372$ đơn vị.
\textbf{Đáp án của câu hỏi là D}
}
\end{ex}

\begin{ex} %[3D2V36]
Một trường tiểu học có 5 khối lớp, mỗi khối có 168 học sinh. Trường đó có tất cả bao nhiêu học sinh?
\choice
{\True $840$ học sinh}
{$800$ học sinh}
{$850$ học sinh}
{$820$ học sinh}
\loigiai{
Tổng số học sinh: $168 \times 5 = 840$ học sinh.
\textbf{Đáp án của câu hỏi là A}
}
\end{ex}

\begin{ex} %[3D2V36]
Bác Năm trồng 4 luống bắp cải, mỗi luống có 125 cây. Đợt rét vừa qua có 35 cây bị chết. Số cây bắp cải còn sống là
\choice
{$500$ cây}
{$475$ cây}
{\True $465$ cây}
{$450$ cây}
\loigiai{
Tổng số cây: $125 \times 4 = 500$ cây. Còn sống: $500 - 35 = 465$ cây.
\textbf{Đáp án của câu hỏi là C}
}
\end{ex}

\textbf{Phần 2. Trắc nghiệm đúng sai}
\setcounterpageref{ex}{0}

\begin{ex} %[3D2V36]
Xét tính đúng sai của phép nhân số có 3 chữ số với số có 1 chữ số:
\choiceTF
{\True Khi nhân số có ba chữ số với số có một chữ số, ta thực hiện nhân lần lượt từ hàng đơn vị sang hàng trăm.}
{$215 \times 3 = 635$.}
{\True $150 \times 4 = 600$.}
{Tích của số có ba chữ số với số có một chữ số không thể vượt quá 999.}
\loigiai{
\begin{itemize}
    \item Ý a) Đúng: Quy tắc tính chuẩn.
    \item Ý b) Sai: Sai vì $215 \times 3 = 645$.
    \item Ý c) Đúng: Tính nhẩm đúng.
    \item Ý d) Sai: Sai vì ví dụ $300 \times 4 = 1200 > 999$.
\end{itemize}
\begin{itemize}
    \item a) \textbf{Đúng}.
    \item b) \textbf{Sai}.
    \item c) \textbf{Đúng}.
    \item d) \textbf{Sai}.
\end{itemize}
}
\end{ex}

\begin{ex} %[3D2V36]
Một xưởng dệt có 3 máy hoạt động, mỗi máy dệt được 230 mét vải mỗi ngày:
\choiceTF
{\True Mỗi ngày cả 3 máy dệt được 690 mét vải.}
{\True Trong 2 ngày một máy dệt được 460 mét vải.}
{Cả xưởng dệt trong 2 ngày dệt được 1200 mét vải.}
{\True Nếu thêm 1 máy nữa cùng công suất thì mỗi ngày dệt được 920 mét vải.}
\loigiai{
\begin{itemize}
    \item Ý a) Đúng: $230 \times 3 = 690\text{ m}$.
    \item Ý b) Đúng: $230 \times 2 = 460\text{ m}$.
    \item Ý c) Sai: Sai vì $690 \times 2 = 1380\text{ m}$.
    \item Ý d) Đúng: $230 \times 4 = 920\text{ m}$.
\end{itemize}
\begin{itemize}
    \item a) \textbf{Đúng}.
    \item b) \textbf{Đúng}.
    \item c) \textbf{Sai}.
    \item d) \textbf{Đúng}.
\end{itemize}
}
\end{ex}

\textbf{Phần 3. Trả lời ngắn}
\setcounterpageref{ex}{0}

\begin{ex} %[3D2V36]
Tính giá trị biểu thức: $148 \times 5 - 240$.
\shortans{500}
\loigiai{
\textbf{Đáp số:} $500$\\
$148 \times 5 = 740$; $740 - 240 = 500$.
}
\end{ex}

\begin{ex} %[3D2V36]
Một mảnh đất hình tam giác có 3 cạnh bằng nhau, mỗi cạnh dài $165\text{ m}$. Tính chu vi mảnh đất đó theo mét.
\shortans{495}
\loigiai{
\textbf{Đáp số:} $495$\\
$165 \times 3 = 495\text{ (m)}$.
}
\end{ex}

\begin{ex} %[3D2V36]
Tìm số tự nhiên $y$, biết: $y : 4 = 175$.
\shortans{700}
\loigiai{
\textbf{Đáp số:} $700$\\
$y = 175 \times 4 = 700$.
}
\end{ex}

\begin{ex} %[3D2V36]
Một rạp chiếu phim có 6 hàng ghế, mỗi hàng có 115 chỗ ngồi. Hỏi rạp có tất cả bao nhiêu chỗ ngồi?
\shortans{690}
\loigiai{
\textbf{Đáp số:} $690$\\
$115 \times 6 = 690$ chỗ ngồi.
}
\end{ex}

\begin{ex} %[3D2V36]
Tính nhanh: $125 \times 2 \times 4$.
\shortans{1000}
\loigiai{
\textbf{Đáp số:} $1000$\\
$125 \times 8 = 1000$.
}
\end{ex}

\begin{ex} %[3D2V36]
Có 3 xe chở gạch, mỗi xe chở 275 viên gạch. Người ta dỡ xuống 350 viên gạch. Hỏi trên các xe còn lại bao nhiêu viên gạch?
\shortans{475}
\loigiai{
\textbf{Đáp số:} $475$\\
Tổng số gạch: $275 \times 3 = 825$ viên. Còn lại: $825 - 350 = 475$ viên.
}
\end{ex}
